\subsection{Bayesian Changepoint Detection}\label{sec:methods_bayesian}

% Che significa "statisticamente indipendenti"? A meno che non ci sia una
% ragione, "idipendenti" è più corto e comprensibile e meno ambiguo. A dire che
% il modello è solido sono già i risultati, non serve sottolinearlo. "evitare",
% è più corto e "escludere ogni possibilità" non mi sembra corretto: esistono
% altre distribuzioni che potrebbero dare risultati diversi, quindi dire "ogni
% possibilità" sembra troppo, è piuttosto un "ridurre le possibilità". Giusto?
% Inoltre, perché solo bias di conferma? Questo migliora la solidità in
% generale.

% In ogni caso, non ho capito come siano stati usati i due modelli: i risultati
% parlano in maniera generale, ma quindi i modelli dicono sempre la stessa cosa,
% su ogni input? Se sì, andrebbe menzionato; se no, c'è qualcosa di sospetto...

% Correggimi se sbaglio, ma quella mi sembra una "spiegazione" di come funziona
% il metodo (simile a quanto si trova online), mentre a noi serve fornire solo i
% metodi usati, le ragioni e i risultati.


%La solidità di questa fase analitica risiede nell'applicazione di due modelli
%bayesiani statisticamente indipendenti, per escludere ogni possibilità di bias
%di conferma. Il primo modello adotta una struttura \textit{Poisson-Gamma}:
%poiché il conteggio delle medaglie è una variabile discreta non negativa, la
%distribuzione di Poisson ne costituisce la descrizione probabilistica naturale,
%mentre la distribuzione Gamma rappresenta il prior coniugato sul parametro di
%intensità $\lambda$. Per ogni possibile punto di rottura temporale $\tau$, il
%modello calcola la log-verosimiglianza marginale della partizione,
%identificando il punto di massima probabilità a posteriori (MAP) come la stima
%della svolta atletica. Il secondo modello analizza il logaritmo del GDP
%nazionale mediante un framework Gaussiano con prior non informativo su media e
%varianza (schema Normal-Inverse-Gamma implicito). 


Per maggiore validità statistica, sono stati applicati due modelli indipendenti:
 per le medagle (variabile discreta non negativa) uno basato su
 \textit{Poisson-Gamma}, per il GDP (logaritmo) un framework Gaussiano con prior
 non informativo su media e varianza (schema Normal-Inverse-Gamma implicito).


% "visualizzata nei grafici come area ombreggiata" messo nei grafici, così non
% contano come parole.


%In entrambi i casi, la ricerca del MAP è vincolata a una finestra temporale di
%$\pm$20 anni rispetto agli anni di hosting, al fine di ancorare l'inferenza e
%ridurre i falsi positivi su serie storiche rumorose. La distribuzione di
%probabilità a posteriori normalizzata fornisce inoltre una misura della
%certezza dell'inferenza: un picco stretto indica un changepoint ben
%identificato, mentre una distribuzione piatta segnala ambiguità nei dati.

La ricerca del changepoint (MAP --- punto di massima probabilità a posteriori) è
stata limitata a una finestra di $\pm$20 anni dall'hosting, per evitare che il
rumore statistico di periodi storici distanti generasse falsi positivi. La
distribuzione di probabilità a posteriori normalizzata fornisce inoltre una
misura della certezza dell'inferenza: un picco stretto indica un changepoint ben
identificato, mentre una distribuzione piatta segnala ambiguità nei dati.